% !TeX program = xelatex
\documentclass[12pt,a4paper]{article}
\usepackage[margin=2cm]{geometry}
\usepackage{fontspec}
\usepackage{amsmath,amssymb}
\usepackage{enumitem}
\setmainfont{Times New Roman}
\setlength{\parindent}{0pt}
\setlength{\parskip}{6pt}
\begin{document}
\begin{center}
{\Large\bfseries Đáp án tuyển sinh vào lớp 10 TP. Hà Nội năm học 2016--2017}\\
{\large Phần Đại số}
\end{center}

\textbf{Câu 1.}
\begin{enumerate}[label=\alph*)]
\item Thay $x=25$:
\[
A=\frac{7}{\sqrt{25}+8}=\frac{7}{13}.
\]
\item Vì $x\ge0$ và $x\ne9$, ta có $\sqrt{x}-3\ne0$ và
\begin{align*}
B&=\frac{\sqrt{x}}{\sqrt{x}-3}+\frac{2\sqrt{x}-24}{(\sqrt{x}-3)(\sqrt{x}+3)}\\
 &=\frac{\sqrt{x}(\sqrt{x}+3)+2\sqrt{x}-24}{(\sqrt{x}-3)(\sqrt{x}+3)}\\
 &=\frac{x+5\sqrt{x}-24}{(\sqrt{x}-3)(\sqrt{x}+3)}
 =\frac{(\sqrt{x}-3)(\sqrt{x}+8)}{(\sqrt{x}-3)(\sqrt{x}+3)}
 =\frac{\sqrt{x}+8}{\sqrt{x}+3}.
\end{align*}
\item Từ câu b),
\[
P=AB=\frac{7}{\sqrt{x}+8}\cdot\frac{\sqrt{x}+8}{\sqrt{x}+3}
 =\frac{7}{\sqrt{x}+3}.
\]
Do $\sqrt{x}+3\ge3$, nên $0<P\le\frac73$. Nếu $P$ nguyên thì $P=1$ hoặc $P=2$.
\[
P=1\Rightarrow \sqrt{x}+3=7\Rightarrow x=16;
\qquad
P=2\Rightarrow \sqrt{x}+3=\frac72\Rightarrow x=\frac14.
\]
Cả hai giá trị đều thỏa điều kiện $x\ge0$, $x\ne9$. Vậy $x=16$ hoặc $x=\frac14$.
\end{enumerate}

\textbf{Câu 2.}

Gọi chiều dài mảnh vườn là $x$ (m), $x>0$. Khi đó chiều rộng là $\dfrac{720}{x}$ (m). Theo đề bài,
\[
(x+10)\left(\frac{720}{x}-6\right)=720.
\]
Nhân hai vế với $x$ và rút gọn:
\[
720x+7200-6x^2-60x=720x
\Longleftrightarrow x^2+10x-1200=0
\Longleftrightarrow (x-30)(x+40)=0.
\]
Vì $x>0$, suy ra $x=30$. Chiều rộng là $720:30=24$ (m). Vậy mảnh vườn dài $30$ m, rộng $24$ m.

\textbf{Câu 3.}
\begin{enumerate}[label=\alph*)]
\item Điều kiện: $x\ne1$, $y\ne-2$. Đặt $u=\dfrac{x}{x-1}$ và $v=\dfrac{1}{y+2}$. Hệ đã cho trở thành
\[
\begin{cases}3u-2v=4,\\2u+v=5.\end{cases}
\]
Giải được $u=2$, $v=1$. Do đó
\[
\frac{x}{x-1}=2\Rightarrow x=2,
\qquad
\frac{1}{y+2}=1\Rightarrow y=-1.
\]
Nghiệm $(x;y)=(2;-1)$ thỏa điều kiện.
\item
\begin{enumerate}[label=(\alph*)]
\item Hoành độ giao điểm thỏa
\[
x^2=3x+m^2-1\Longleftrightarrow x^2-3x-m^2+1=0.
\]
Biệt thức
\[
\Delta=9-4(-m^2+1)=4m^2+5>0
\]
với mọi $m$. Phương trình có hai nghiệm phân biệt, nên $(d)$ luôn cắt $(P)$ tại hai điểm phân biệt.
\item Theo Viète, $x_1+x_2=3$ và $x_1x_2=1-m^2$. Điều kiện đề bài tương đương
\begin{align*}
(x_1+1)(x_2+1)=1
&\Longleftrightarrow x_1x_2+x_1+x_2+1=1\\
&\Longleftrightarrow (1-m^2)+3+1=1\\
&\Longleftrightarrow m^2=4\\
&\Longleftrightarrow m=\pm2.
\end{align*}
Vậy $m=2$ hoặc $m=-2$.
\end{enumerate}
\end{enumerate}

\textbf{Câu 5.}

Điều kiện xác định: $x\ge-6$, $y\ge-6$. Từ giả thiết,
\[
x-\sqrt{x+6}=\sqrt{y+6}-y
\Longleftrightarrow x+y=\sqrt{x+6}+\sqrt{y+6}.
\]
Đặt $S=x+y$. Vế phải không âm nên $S\ge0$. Bất đẳng thức $(\sqrt a+\sqrt b)^2\le2(a+b)$ với $a,b\ge0$ cho
\[
S^2\le2\bigl((x+6)+(y+6)\bigr)=2S+24,
\]
suy ra $-4\le S\le6$. Mặt khác, bình phương đẳng thức $S=\sqrt{x+6}+\sqrt{y+6}$ cho
\[
S^2=S+12+2\sqrt{(x+6)(y+6)},
\]
nên
\[
(S+3)(S-4)=S^2-S-12=2\sqrt{(x+6)(y+6)}\ge0.
\]
Kết hợp với $S\ge0$ suy ra $S\ge4$. Vậy $4\le P=S\le6$.

Nếu $S=4$, đẳng thức trên buộc $(x+6)(y+6)=0$. Kết hợp $x+y=4$, ta được $(x,y)=(-6,10)$ hoặc $(10,-6)$; do đó giá trị nhỏ nhất là $4$.

Nếu $S=6$, dấu bằng trong bất đẳng thức $(\sqrt a+\sqrt b)^2\le2(a+b)$ xảy ra khi $a=b$, tức $x+6=y+6$. Cùng với $x+y=6$, suy ra $x=y=3$; do đó giá trị lớn nhất là $6$.

\end{document}
